\subsection{交换引理}

引理 1. 设 \(\mathbf{A}\) 是一个交换环， \(\mathbf{E}\) 是一个 \(\mathbf{A}\) -代数， \((x_{i})_{1\leqslant i\leqslant n}\) 是 \(\mathbf{E}\) 的元素组成的一个有限序列， \((\lambda_i)_{1\leqslant i\leqslant n}\) 是 \(\mathbf{A}\) 的元素组成的一个有限序列， \(y\) 是 \(\mathbf{E}\) 的一个元素；假设

\[
x _ {i} y = \lambda_ {i} y x _ {i} \quad \text { for } 1 \leqslant i \leqslant n.\tag{13}
\]

那么

\[
\left(x _ {1} x _ {2} \dots x _ {n}\right) y = \left(\lambda_ {1} \lambda_ {2} \dots \lambda_ {n}\right) y \left(x _ {1} x _ {2} \dots x _ {n}\right).\tag{14}
\]

该引理对 \(n = 1\) 是平凡的，我们对 \(n \geqslant 2\) 作归纳论证。现在

\[
\begin{array}{r l} (x _ {1} x _ {2} \dots x _ {n}) y & = (x _ {1} \dots x _ {n - 1}) (x _ {n} y) \\ & = (x _ {1} \dots x _ {n - 1}) (\lambda_ {n} y x _ {n}) = \lambda_ {n} ((x _ {1} \dots x _ {n - 1}) y) x _ {n}, \end{array}
\]

根据归纳假设，其等于

\[
\lambda_ {n} \left(\lambda_ {1} \dots \lambda_ {n - 1}\right) y \left(x _ {1} \dots x _ {n - 1}\right) x _ {n} = \left(\lambda_ {1} \dots \lambda_ {n - 1} \lambda_ {n}\right) y \left(x _ {1} \dots x _ {n - 1} x _ {n}\right),
\]

由此即得引理。

引理 2. 设 A 是一个交换环，E 是一个 A-代数， \((x_{i})_{1\leqslant i\leqslant n}\) 和 \((y_{i})_{1\leqslant i\leqslant n}\) 是 E 的由 \(n\) 个元素组成的两个有限序列；假设对 \(1\leqslant j\leqslant i\leqslant n\) ，

\[
x _ {i} y _ {j} = \lambda_ {i j} y _ {j} x _ {i} \quad \text { with } \lambda_ {i j} \in A.\tag{15}
\]

那么

\[
\left(x _ {1} x _ {2} \dots x _ {n}\right) \left(y _ {1} y _ {2} \dots y _ {n}\right) = \left(\prod_ {i > j} \lambda_ {i j}\right) \left(x _ {1} y _ {1}\right) \left(x _ {2} y _ {2}\right) \dots \left(x _ {n} y _ {n}\right).\tag{16}
\]

该引理对 \(n = 1\) 是平凡的，我们对 \(n\) （ \(n \geqslant 2\) ）再次作归纳论证。根据引理 1，

\[
\begin{array}{r l} (x _ {1} \dots x _ {n}) (y _ {1} \dots y _ {n}) & = x _ {1} (x _ {2} \dots x _ {n}) y _ {1} (y _ {2} \dots y _ {n}) \\ & = \left(\prod_ {i > 1} \lambda_ {i 1}\right) (x _ {1} y _ {1}) (x _ {2} \dots x _ {n}) (y _ {2} \dots y _ {n}) \end{array}
\]

然后只需应用归纳假设即可得到 (16)。

对元素的每个族 \(\lambda = (\lambda_{ij})\) （元素属于 A，且 \(1 \leqslant j < i \leqslant n\) ），以及对每个置换 \(\sigma \in \mathfrak{S}_n\) ，我们记

\[
\varepsilon_ {\sigma} (\lambda) = \prod_ {i > j, \sigma^ {- 1} (i) < \sigma^ {- 1} (j)} \lambda_ {i j} = \prod_ {i <   j, \sigma (i) > \sigma (j)} \lambda_ {\sigma (i), \sigma (j)}.\tag{17}
\]

注意，当 \(\mathbf{A} = \mathbf{Z}\) 且 \(\lambda_{ij} = -1\) 对每个有序对 \((i,j)\) （满足 \(1 \leqslant j < i \leqslant n\) ）成立时， \(\varepsilon_{\sigma}(\lambda)\) 就是符号 \(\varepsilon_{\sigma}\) （置换 \(\sigma\) 的符号）（I, §5, 第 7 号）。

引理 3. 设 A 是一个交换环，E 是一个 A-代数， \((x_{i})_{1\leqslant i\leqslant n}\) 是 E 的元素组成的一个有限序列，并假设对每个整数有序对 \((i,j)\) （满足 \(1\leqslant j < i\leqslant n\) ），

\[
x _ {i} x _ {j} = \lambda_ {i j} x _ {j} x _ {i} \quad \text { with } \lambda_ {i j} \in A.\tag{18}
\]

那么，对每个置换 \(\sigma \in \mathfrak{S}_n\) ，

\[
x _ {\sigma (1)} x _ {\sigma (2)} \dots x _ {\sigma (n)} = \varepsilon_ {\sigma} (\lambda) x _ {1} x _ {2} \dots x _ {n}.\tag{19}
\]

该引理对 \(n = 1\) 和 \(n = 2\) 是平凡的；我们对 \(n\) （ \(n \geqslant 3\) ）作归纳。若 \(\sigma(n) = n\) ，则关系式 (19) 由归纳假设得出。因此设 \(\sigma(n) = k, k \neq n\) ，并令 \(\tau\) 为 \([1, n]\) 上由下式定义的置换

\[
\left\{ \begin{array}{l l} \tau (i) = i & \text { for } i <   k \\ \tau (i) = i + 1 & \text { for } k \leqslant i <   n \\ \tau (n) = k. \end{array} \right.
\]

设 \(\pi = \tau^{-1} \circ \sigma\) ；置换 \(\pi\) 保持 \(n\) 不动；现在 \(\sigma = \tau \circ \pi\) ，因此，记 \(y_i = x_{\tau(i)}\) ，就有 \(y_{\pi(i)} = x_{\sigma(i)}\) 。若 \(i \neq n\) 且 \(j \neq n\) ，关系 \(\pi(i) > \pi(j)\) 与 \(\sigma(i) > \sigma(j)\) 等价（因为 \(\tau\) 是从 \([1, n - 1]\) 到 \([1, n]\) 的严格递增映射。对 \(i \neq n\) , \(j \neq n\) 和 \(\sigma(i) > \sigma(j)\) ，

\[
y _ {\pi (i)} y _ {\pi (j)} = x _ {\sigma (i)} x _ {\sigma (j)} = \lambda_ {\sigma (i), \sigma (j)} x _ {\sigma (j)} x _ {\sigma (i)} = \lambda_ {\sigma (i), \sigma (j)} y _ {\pi (j)} y _ {\pi (i)}
\]

因此，根据归纳假设（利用 \(\pi(n) = n\) ）:

\[
y _ {\pi (1)} y _ {\pi (2)} \dots y _ {\pi (n)} = \left(\prod_ {i <   j <   n, \sigma (i) > \sigma (j)} \lambda_ {\sigma (i), \sigma (j)}\right) y _ {1} y _ {2} \dots y _ {n}
\]

即

\[
x _ {\sigma (1)} x _ {\sigma (2)} \dots x _ {\sigma (n)} = \left(\prod_ {i <   j <   n, \sigma (i) > \sigma (j)} \lambda_ {\sigma (i), \sigma (j)}\right) x _ {\tau (1)} \dots x _ {\tau (n)}.\tag{20}
\]

现在

\[
x _ {\tau (1)} \dots x _ {\tau (n)} = x _ {1} \dots x _ {k - 1} x _ {k + 1} \dots x _ {n} x _ {k}
\]

而由引理 1，这等于

\[
\left(\prod_ {j > k} \lambda_ {j k}\right) x _ {1} \dots x _ {n} = \left(\prod_ {\sigma (i) > \sigma (n)} \lambda_ {\sigma (i), \sigma (n)}\right) x _ {1} \dots x _ {n}.\tag{21}
\]

最后，(20) 与 (21) 给出

\[
x _ {\sigma (1)} \dots x _ {\sigma (n)} = \alpha . x _ {1} \dots x _ {n}
\]

其中

\[
\begin{array}{l} \alpha = \left(\prod_ {i <   j <   n, \sigma (i) > \sigma (j)} \lambda_ {\sigma (i), \sigma (j)}\right) \cdot \left(\prod_ {i <   n, \sigma (i) > \sigma (n)} \lambda_ {\sigma (i), \sigma (n)}\right) \\ = \prod_ {i <   j, \sigma (i) > \sigma (j)} \lambda_ {\sigma (i), \sigma (j)} = \varepsilon_ {\sigma} (\lambda) \end{array}
\]

这就完成了引理 3 的证明。

